% Generated by `taskgen course examples` from crypto.rsa_factor_from_d (seed 4).
% Do not edit: change examples.yaml and regenerate.
\begin{example}[разложение $N$ по ключам $e$ и $d$]\label{ex:rsa-from-d}
Известны открытый ключ RSA $(N, e) = (1073, 5)$ и соответствующий закрытый ключ $d = 605$. Разложить $N$ на множители вероятностным алгоритмом: $ed - 1 = 2^t r$, $r$ нечётно; для оснований $a = 2, 3, 5, \ldots$ (по порядку) искать нетривиальный квадратный корень из единицы среди $a^{r}, a^{2r}, \ldots, a^{2^t r} \bmod N$.

\par\noindent\textit{Решение.} $ed - 1 = 5 \cdot 605 - 1 = 3024 = 2^{4} \cdot 189$, т.\,е. $t = 4$, $r = 189$. Так как $\lambda(N) \mid ed - 1$, для $a$, взаимно простого с $N$, $a^{2^t r} \equiv 1 \pmod{N}$. Вычисляем $x_0 = a^r \bmod N$ (бинарным алгоритмом) и возводим в квадрат, пока не получим $1$. Если $x_0 = 1$ или перед единицей стоит $-1 \equiv 1072$, основание неудачно; иначе $x^2 \equiv 1$, $x \not\equiv \pm 1$, т.\,е. $N \mid (x - 1)(x + 1)$, но $N$ не делит ни один множитель, и $\gcd(x - 1, N)$ --- нетривиальный делитель.\par\noindent $a = 2$: $x_{0} = 771$, $x_{1} = 1072$, $x_{2} = 1$. Перед единицей стоит $-1$ --- неудача, берём следующее основание.\par\noindent $a = 3$: $x_{0} = 887$, $x_{1} = 260$, $x_{2} = 1$. $x = 260$ --- нетривиальный корень из $1$: $\gcd(260 - 1, 1073) = 37$.\par $N / 37 = 29$.\[N = 29 \cdot 37 = 1073.\]
\end{example}
